\section{Relation to other perspectives}
\label{sec:learning_mech_and_mechinterp}

There are several ongoing approaches to developing explanatory scientific theory of deep learning, each adopting a different perspective and using different sets of tools.
We believe that these perspectives are essentially all complementary: all either \textit{directly seek} a mechanics of learning or would \textit{symbiotically benefit} from one.

\paragraph{The statistical perspective.}
The rich tradition of classical learning theory remains influential today.\footnote{
The Simons Institute for the Theory of Computing has provided an important substrate for developing this statistical perspective of deep learning through \href{https://deepfoundations.ai/}{collaborations} and \href{https://simons.berkeley.edu/workshops/deep-learning-theory}{seminars}.
}
\citet{bartlett:2021-statistical-viewpoint} offer a lucid summary of its central framing: any statistical prediction method must balance \textit{expressivity} (to represent the richness of real data), \textit{complexity control} (to make the most of finite training data), and \textit{computational efficiency} (to yield practical algorithms).
It is apparent that deep learning is sufficiently expressive, but it is not clear how a good function is selected from this enormous function class, nor why simple gradient methods suffice to train such complex beasts.
The modern statistical viewpoint suggests two answers: deep learning has an \textit{implicit inductive bias} towards simple, well-generalizing functions \citep{wilsonposition}, and despite their nonconvexity, the \textit{very high dimensionality} (overparameterization) of neural networks makes optimization easy.


These questions are good ones, and we believe these answers are basically correct.
The challenge is now to make them precise in the case of neural networks.
It is clear that doing so will require taking a close look at the nature of the training process.
Only once we have done so will we be able to back out \textit{how} this implicit bias arises and \textit{why} gradient methods suffice for optimization.
We do not believe these answers will be generic statements, but instead critically rely on important properties of deep learning and of natural data.
The statistical perspective thus leads naturally to a serious scientific study of the mechanics of training.

\paragraph{The information-theoretic perspective.}
A closely-related approach seeks to explain deep learning in terms of information-theoretic ideas.
In this view, learning is a process of extracting information from datasets, and a learning system works when it extracts information useful for prediction while discarding irrelevant information.
This perspective hopes to understand learning as \textit{compression} of the dataset into either the model's parameters or its hidden representations, with good generalization resulting when this compression is successful \citep{shwartz2017opening,xu2017information}.

We find this perspective insightful, and it seems likely to us that a picture of this nature will hold.
As with the statistical perspective, a major remaining question is how to make this view concrete and actionable: how do the architecture and training process of deep learning interact to actually \textit{implement} this compression, and what factors make it more or less successful?
Doing this, too, will require taking a close look at the nature of the training process, the architecture, the data, and their interactions.
The information-theoretic perspective thus also leads naturally to a serious scientific study of the mechanics of training.

\paragraph{Physics of deep learning.}
This community descends from the older physics of machine learning lineage \citep{hopfield1982neural,amit1985spin,gardner1988space} and essentially seeks satisfying average-case theories of neural network learning \citep{zdeborova2020understanding,bahri2020statistical,michaudphysics,ringel2025applications}.\footnote{In the modern day, this community is mediated in part by recurring events at the \href{https://www.kitp.ucsb.edu/activities/deeplearning23}{Kavli Institute for Theoretical Physics}, the \href{https://aspenphys.org/event/theoretical-physics-for-artificial-intelligence}{Aspen Center for Theoretical Physics}, and the \href{https://leshouches2022.github.io/}{Les Houches School of Physics}, and organizations such as the \href{https://iaifi.org/}{NSF AI Institute for Artificial Intelligence and Fundamental Interactions} and the \href{https://www.physicsoflearning.org/}{Simons collaboration on the physics of learning and neural computation}.}
The close relationship between physics and machine learning was recognized by the \href{https://www.nobelprize.org/prizes/physics/2024/summary/}{2024 Nobel Prize in Physics}.
This approach is in line with (and has largely shaped) the perspective presented in this paper, and the project of this community is arguably the development of a mechanics of learning.
The challenge then is to clarify important problems and coordinate effort for efficient progress.

\paragraph{Perspectives from neuroscience.}
Several approaches to developing a science of the brain suggest approaches to developing a science of deep learning.
One approach starts from hypotheses about neural systems --- for example, that their computation amounts to some form of approximate probabilistic inference --- and seeks to make deductions and predictions from this hypothesis \citep{dayan1995helmholtz,friston2010free}.
Some of these predictions seem to hold suspiciously well in deep learning: see, for example, the case of edge-selective cells in the visual cortex \citep{olshausen1996emergence} and edge-selective receptive fields in convolutional networks \citep[e.g.][]{zeiler2014visualizing}.
Another approach termed \textit{systems neuroscience} seeks to directly decompose subsets of the brain into interpretable circuits and reverse-engineer the structure of their learned representations \citep{chung2021neural,bernardi2020geometry,kriegeskorte2008representational}.
This approach resembles mechanistic interpretability, which has adopted some of its methods and intuitions.

We expect and encourage this dialogue to continue, and it seems plausible that some of these high-level hypotheses about the brain --- e.g., that the brain admits at least a partial decomposition into interpretable circuits and that local circuits implicitly solve inference tasks --- will turn out to be true of deep learning.
The reasons these facts are true, if indeed they are, is surely bound up in the dynamical way learning actually happens.
A study of the mechanics of learning is thus important to the continued exploration of these ideas.

\paragraph{Developmental interpretability/singular learning theory.}
This approach, which grew out of the mechanistic interpretability community, seeks first-principles predictive theories of neural network learning based on the singular learning theory framework of \citet{watanabe2009algebraic}, emphasizing a Bayesian perspective and aiming to understand training as a process of sequential phase transitions mediated by the geometry of the loss landscape \citep{hoogland2023developmental}.
We see this community as seeking the same goal we suggest here --- a fundamental mechanics of learning, and a rigorous foundation for interpretability --- but with a toolkit that differs from the other listed perspectives.
There is potential for fruitful cross-pollination and tool-sharing between these different approaches.

\paragraph{Science of deep learning.}
It has long been appreciated by practitioners that machine learning is largely a practice of trial and error and that it may be possible and beneficial to systematize it \citep{langley1988machine,gal2015science,rahimi2017nips_test_of_time,baraniuk2020science}.
Indeed, much of the rapid empirical progress of the last decade resulted from systematic organization around agreed-upon benchmark tasks \citep{donoho2024data}.
Nonetheless, the training and application of large models remains more alchemy than science.
We believe that a fundamental mechanics of the learning process is the foundation on which this science will finally be built.






\subsection{\texorpdfstring{Learning mechanics $\rightleftarrows$ mechanistic interpretability}{Learning mechanics <-> mechanistic interpretability}}
\label{subsec:learning_mech_and_mechinterp}

We discuss mechanistic interpretability specially because there is a unique opportunity for cooperation.
Mechanistic interpretability aims to understand trained neural networks by identifying the internal mechanisms
--- features, circuits, and learned algorithms ---
that give rise to their behavior. %
At its core, this approach is guided by the belief that neural networks admit a  human-understandable, mechanistic description that can be uncovered through careful empirical reverse engineering.\footnote{The mechanistic interpretability community does not yet share a formal definition of what constitutes a ``mechanistic description,'' though see~\citet{geiger2025causal} for a recently proposed causal framing. Informally, many researchers proceed under a set of working assumptions: (1) that neural networks encode the state of internal computational variables in their activations, often referred to as ``features''; (2) that successive layers transform and combine these features in structured ``circuits''; and (3) that, taken together, these circuits implement algorithms that admit some level of human-understandable description.}
This approach has already borne fruit: many visually striking or interpretable mechanisms have been discovered in large models to date \citep{olah2020zoom,templeton2024scaling,engels2024not,gurnee2025when,lindsey2025biology}.%
\footnote{Mechanistic interpretability is deeply associated with \href{https://anthropic.com}{Anthropic} and the \href{https://www.alignmentforum.org/}{AI safety} and Effective Altruism communities, though it is increasingly pursued in academic labs. We also note that mechanistic interpretability has recently split into an \textit{ambitious} camp that hopes to develop full, explanatory scientific theory and a \textit{pragmatic} camp which is mostly interested in targeted interventions for particular cases. See also~\mbox{\cite{saphra2024mechanistic}} for a discussion of the origin of the term ``mechanistic interpretability'' and the dynamics of its community.}

This is a complementary perspective to our own and presents a wonderful opportunity for symbiosis.
At time of writing, mechanistic interpretability remains largely a qualitative science, more reliant on human-judged empirics than on compact mathematical principles or simple governing laws.
This is quite natural: semantically-meaningful functions resist mathematical characterization.%
\footnote{For example, try writing down a function that can classify dogs vs. cats from image pixel values. The difficulty of expressing such functions in mathematics is why we invented deep learning in the first place!}
On the other hand, a mechanics of learning would be quantitative by definition, but by the same token will be too low-level to answer important questions of semantic meaning on its own.
These approaches study the same system --- i.e., deep learning --- at different levels of abstraction, and so of course they can (and should) work together for mutual gain.
Calls for rigorous foundations for interpretability have been steadily growing \citep{sharkey2025open,joshi2026causality,greenspan2026towards}, and this is one thing learning mechanics can and should seek to help provide.
In turn, mechanistic interpretability offers learning mechanics a rich and growing tableau of empirical phenomena ripe for the development of explanatory mathematical theory.












\paragraph{Learning mechanics $\rightarrow$ mechanistic interpretability.}
We emphasize two complementary avenues through which learning mechanics can support mechanistic interpretability: formalizing core assumptions and explaining how mechanisms develop through training.


\emph{Formalizing core assumptions.} Learning mechanics can make explicit, formalize, and, where necessary, challenge the core and often implicit assumptions that guide interpretability research.
These include:
\begin{itemize}
  \item \emph{linear representability} --- that features correspond to meaningful directions in activation space \citep{mikolov2013linguistic, park2023linear, nanda2023emergent, marks2023geometry, jiang2024origins, csordas2024recurrent};
  \item \emph{locality} --- features and circuits are localizable to particular subsets of model components\citep{meng2022locating, wang2022interpretability, conmy2023towards, arora2025language};
  \item \emph{sparsity} --- that individual features and circuits are activated or functionally relevant on only a small fraction of inputs ~\citep{cunningham2023sparse, bricken2023monosemanticity}; and
  \item \emph{compositionality} --- that complex network representations and computations arise from the composition of simpler, modular sub-mechanisms \citep{thorpe1989local,smolensky1990tensor,lepori2023break,schug2023discovering,ramesh2023compositional}.
\end{itemize}
These core assumptions underpin the identification, isolation, and analysis of the internal mechanisms of trained neural networks in mechanistic interpretability research.
A mathematical theory of learning offers a way to clarify the regimes in which these assumptions hold, the conditions under which they fail, and the sense in which they can be derived from training dynamics and data statistics (see \cref{openquestion:what_are_features}).



\emph{Explaining how mechanisms develop through training.} Mechanistic interpretability has generally prioritized describing \emph{what} mechanisms trained neural networks have learned,
and there remains a rich opportunity for work which aims to explain \emph{how} and \emph{why} such mechanisms form in the first place.
There is already substantial interest within parts of the interpretability community in this dynamical/theoretical perspective, including work on the formation of induction heads~\citep{elhage2021mathematical, olsson2022context}, grokking and progress measures~\citep{nanda2023progress}, sudden phase transitions in circuit formation \citep{elhage2022toy, chen2023sudden, gopalani2024abrupt,park2024emergence}, and the research program of developmental interpretability~\citep{hoogland2023developmental,hoogland2025loss}, discussed earlier.
Our goal is not to replace these efforts but to encourage deeper engagement between mechanistic interpretability and the broader landscape of mathematically grounded ideas and tools in learning mechanics.
Echoing~\cite{saphra2022creationism}, we hope that
learning mechanics can play a role analogous to evolution in biology: just as \emph{``nothing in biology makes sense except in the light of evolution,''} the internal mechanisms of trained networks may be most naturally understood in the light of the processes that give rise to them.


\paragraph{Learning mechanics $\leftarrow$ mechanistic interpretability.}
Conversely, learning mechanics has been deeply influenced by the empirical discoveries of mechanistic interpretability, which often identify concrete phenomena that invite first-principles explanation.
Mechanistic interpretability places the structure of data at the center of its analyses, revealing settings in which the relationship between input structure and learned mechanisms is especially clear~\citep{nanda2023progress, shai2024transformers}.
By contrast, much of classical deep learning theory has relied on highly simplified data models, leaving a gap between theoretical predictions and behaviors observed in practice.
In this way, mechanistic interpretability helps bridge this gap by providing learning mechanics with concrete, well-defined targets for theoretical modeling.

Several such observations have already proven influential in stimulating work in learning mechanics, including the emergence of induction heads for in-context learning \citep{bietti2023birth, reddy2023mechanistic,nichani2024transformers}, the role of Fourier features in algebraic tasks \citep{morwani2023feature,kunin2025alternating,marchetti2026sequential}, and the geometry of features arising from the structure of correlations in the data \citep{engels2024not, prieto2025correlations, karkada2026symmetry}.
Just as the development of physics was often driven by empirical discoveries in adjacent fields, we expect progress in learning mechanics to be driven by theorists who take seriously empirical phenomena, including those uncovered by the mechanistic interpretability community, and seek to explain them.
